\chapter{系统分析与设计}

\section{需求分析}

校园网建设需要满足多区域接入、业务隔离、稳定转发和安全访问控制等要求。根据使用场景，可将网络划分为教学区、办公区、宿舍区、服务器区和公共无线区。各区域通过虚拟局域网进行逻辑隔离，并在核心层统一完成路由转发和策略控制。

\subsection{系统概要}

系统采用核心层、汇聚层和接入层的三层网络架构。核心层负责高速转发与冗余互联，汇聚层负责区域流量汇聚与策略控制，接入层负责终端接入与基础安全限制。整体结构如图~\ref{fig:topology} 所示。

\begin{figure}[htbp]
  \centering
  \fbox{%
    \begin{minipage}[c][5.2cm][c]{0.78\textwidth}
      \centering
      核心交换区\\[0.5em]
      $\downarrow$\\[0.5em]
      教学区 \quad 办公区 \quad 宿舍区 \quad 服务器区 \quad 无线区
    \end{minipage}
  }
  \caption{校园网总体拓扑示意图}
  \label{fig:topology}
\end{figure}

\subsection{系统流程}

用户终端接入网络后，首先由接入层交换机识别所属区域，再根据虚拟局域网配置进入对应逻辑网络。跨区域访问由汇聚层和核心层完成三层转发，并通过访问控制策略限制非授权访问。服务器区对外提供统一业务入口，减少内部系统暴露面。

\section{概要设计}

概要设计阶段主要确定网络层级、设备角色、地址规划和安全边界。为便于后续维护，本文采用按区域划分的地址规划方案，并为教学、办公、宿舍和服务器等业务分别预留独立网段。

\begin{table}[htbp]
  \centering
  \caption{校园网主要区域规划}
  \label{tab:area-plan}
  \begin{tabular}{lll}
    \toprule
    区域 & 主要用途 & 设计重点 \\
    \midrule
    教学区 & 多媒体教室、实验室 & 稳定接入、实验隔离 \\
    办公区 & 行政办公、教师办公 & 访问控制、数据安全 \\
    宿舍区 & 学生终端接入 & 高并发、带宽管理 \\
    服务器区 & 教务、资源与认证系统 & 安全边界、冗余保护 \\
    \bottomrule
  \end{tabular}
\end{table}

\subsection{业务流程}

业务访问流程可以概括为终端接入、身份识别、地址分配、策略匹配和资源访问五个阶段。对于敏感业务，系统应限制访问来源，并记录关键访问日志。

\subsection{功能模块介绍}

本方案的主要模块包括网络拓扑设计、地址规划、路由配置、安全策略和运维管理。各模块之间相互关联，其中地址规划和安全策略会直接影响后续配置复杂度。

